\subsection{CodeFeedback selection and evaluation}\label{app:coding}
We use the Python subset of CodeFeedback \citep{zheng2024opencodeinterpreter}, distributed for PiSSA \citep{meng2024pissa}. After exact prompt deduplication, a fixed split contains 20,000 training and 2,000 development examples. Exact prompt matching against the test sets removes no additional examples. All methods use one epoch, effective batch size 32, maximum training length 1,024, a cosine schedule with 3\% warmup, and the same attention/MLP projections. LoRA, DoRA, PiSSA and MiLoRA train at rank 7. OFT uses block size 32 and HRA uses 16 Householder vectors. PiSSA/MiLoRA's converted rank-14 difference adapters retain the same training capacity.

For each method/backbone, the mean development response NLL across seeds 13, 37 and 73 selects one of $\{10^{-5},3\times10^{-5},10^{-4},3\times10^{-4},10^{-3}\}$, with ties favoring the lower LR. NLL pools valid response-token losses, with a 4,096-token evaluation limit. Thus the 180 development runs yield 36 selected test checkpoints, each using the best LR within the tested grid. 

Test evaluation uses EvalPlus base completion prompts, greedy decoding, and one completion per problem, capped at 1,024 new tokens. We evaluate 164 HumanEval \citep{chen2021humaneval} and 378 MBPP \citep{austin2021mbpp} problems. Original and strengthened tests grade the same completion. \Cref{tab:coding-full} reports both versions. The base model has one fixed evaluation, so its row has no training-seed SD.

General retention uses the fixed FineWiki-200 corpus, a 768-token limit, and valid-token-pooled NLL with float32 logits. Each method's forgetting is its NLL minus the corresponding unadapted model's NLL. 

\begin{table}[!htbp]
\caption{CodeFeedback results for all original and strengthened-test pass@1 scores (\%). A single greedy completion is graded under both test suites. LR minimizes mean held-out response NLL. Only selected checkpoints have test evaluations. The Base row has no training-seed SD.}
\label{tab:coding-full}
\begin{center}
\setlength{\tabcolsep}{3pt}
\begin{tabular*}{\linewidth}{@{\extracolsep{\fill}}llccccc@{}}\toprule
Model & Method & Selected LR & HumanEval & HE+ & MBPP & MBPP+ \\ \midrule
Qwen7B & Base & --- & $\result{ext_code_qwen_base_he}$ & $\result{ext_code_qwen_base_hep}$ & $\result{ext_code_qwen_base_mb}$ & $\result{ext_code_qwen_base_mbp}$ \\
 & \textcolor{LoRAColor}{LoRA} & $1\times10^{-4}$ & $\result{ext_code_qwen_lora_0p0001_he_mean}_{\pm\result{ext_code_qwen_lora_0p0001_he_sd}}$ & $\result{ext_code_qwen_lora_0p0001_hep_mean}_{\pm\result{ext_code_qwen_lora_0p0001_hep_sd}}$ & $\result{ext_code_qwen_lora_0p0001_mb_mean}_{\pm\result{ext_code_qwen_lora_0p0001_mb_sd}}$ & $\result{ext_code_qwen_lora_0p0001_mbp_mean}_{\pm\result{ext_code_qwen_lora_0p0001_mbp_sd}}$ \\
 & \textcolor{OFTColor}{OFT} & $1\times10^{-4}$ & $\result{ext_code_qwen_oft_0p0001_he_mean}_{\pm\result{ext_code_qwen_oft_0p0001_he_sd}}$ & $\result{ext_code_qwen_oft_0p0001_hep_mean}_{\pm\result{ext_code_qwen_oft_0p0001_hep_sd}}$ & $\result{ext_code_qwen_oft_0p0001_mb_mean}_{\pm\result{ext_code_qwen_oft_0p0001_mb_sd}}$ & $\result{ext_code_qwen_oft_0p0001_mbp_mean}_{\pm\result{ext_code_qwen_oft_0p0001_mbp_sd}}$ \\
 & \textcolor{DoRAColor}{DoRA} & $1\times10^{-4}$ & $\result{ext_code_qwen_dora_0p0001_he_mean}_{\pm\result{ext_code_qwen_dora_0p0001_he_sd}}$ & $\result{ext_code_qwen_dora_0p0001_hep_mean}_{\pm\result{ext_code_qwen_dora_0p0001_hep_sd}}$ & $\result{ext_code_qwen_dora_0p0001_mb_mean}_{\pm\result{ext_code_qwen_dora_0p0001_mb_sd}}$ & $\result{ext_code_qwen_dora_0p0001_mbp_mean}_{\pm\result{ext_code_qwen_dora_0p0001_mbp_sd}}$ \\
 & \textcolor{PiSSAColor}{PiSSA} & $1\times10^{-5}$ & $\result{ext_code_qwen_pissa_1em05_he_mean}_{\pm\result{ext_code_qwen_pissa_1em05_he_sd}}$ & $\result{ext_code_qwen_pissa_1em05_hep_mean}_{\pm\result{ext_code_qwen_pissa_1em05_hep_sd}}$ & $\result{ext_code_qwen_pissa_1em05_mb_mean}_{\pm\result{ext_code_qwen_pissa_1em05_mb_sd}}$ & $\result{ext_code_qwen_pissa_1em05_mbp_mean}_{\pm\result{ext_code_qwen_pissa_1em05_mbp_sd}}$ \\
 & \textcolor{MiLoRAColor}{MiLoRA} & $1\times10^{-4}$ & $\result{ext_code_qwen_milora_0p0001_he_mean}_{\pm\result{ext_code_qwen_milora_0p0001_he_sd}}$ & $\result{ext_code_qwen_milora_0p0001_hep_mean}_{\pm\result{ext_code_qwen_milora_0p0001_hep_sd}}$ & $\result{ext_code_qwen_milora_0p0001_mb_mean}_{\pm\result{ext_code_qwen_milora_0p0001_mb_sd}}$ & $\result{ext_code_qwen_milora_0p0001_mbp_mean}_{\pm\result{ext_code_qwen_milora_0p0001_mbp_sd}}$ \\
 & \textcolor{HRAColor}{HRA} & $1\times10^{-3}$ & $\result{ext_code_qwen_hra_0p001_he_mean}_{\pm\result{ext_code_qwen_hra_0p001_he_sd}}$ & $\result{ext_code_qwen_hra_0p001_hep_mean}_{\pm\result{ext_code_qwen_hra_0p001_hep_sd}}$ & $\result{ext_code_qwen_hra_0p001_mb_mean}_{\pm\result{ext_code_qwen_hra_0p001_mb_sd}}$ & $\result{ext_code_qwen_hra_0p001_mbp_mean}_{\pm\result{ext_code_qwen_hra_0p001_mbp_sd}}$ \\
\addlinespace[3pt]
Llama8B & Base & --- & $\result{ext_code_llama_base_he}$ & $\result{ext_code_llama_base_hep}$ & $\result{ext_code_llama_base_mb}$ & $\result{ext_code_llama_base_mbp}$ \\
 & \textcolor{LoRAColor}{LoRA} & $3\times10^{-4}$ & $\result{ext_code_llama_lora_0p0003_he_mean}_{\pm\result{ext_code_llama_lora_0p0003_he_sd}}$ & $\result{ext_code_llama_lora_0p0003_hep_mean}_{\pm\result{ext_code_llama_lora_0p0003_hep_sd}}$ & $\result{ext_code_llama_lora_0p0003_mb_mean}_{\pm\result{ext_code_llama_lora_0p0003_mb_sd}}$ & $\result{ext_code_llama_lora_0p0003_mbp_mean}_{\pm\result{ext_code_llama_lora_0p0003_mbp_sd}}$ \\
 & \textcolor{OFTColor}{OFT} & $3\times10^{-4}$ & $\result{ext_code_llama_oft_0p0003_he_mean}_{\pm\result{ext_code_llama_oft_0p0003_he_sd}}$ & $\result{ext_code_llama_oft_0p0003_hep_mean}_{\pm\result{ext_code_llama_oft_0p0003_hep_sd}}$ & $\result{ext_code_llama_oft_0p0003_mb_mean}_{\pm\result{ext_code_llama_oft_0p0003_mb_sd}}$ & $\result{ext_code_llama_oft_0p0003_mbp_mean}_{\pm\result{ext_code_llama_oft_0p0003_mbp_sd}}$ \\
 & \textcolor{DoRAColor}{DoRA} & $3\times10^{-4}$ & $\result{ext_code_llama_dora_0p0003_he_mean}_{\pm\result{ext_code_llama_dora_0p0003_he_sd}}$ & $\result{ext_code_llama_dora_0p0003_hep_mean}_{\pm\result{ext_code_llama_dora_0p0003_hep_sd}}$ & $\result{ext_code_llama_dora_0p0003_mb_mean}_{\pm\result{ext_code_llama_dora_0p0003_mb_sd}}$ & $\result{ext_code_llama_dora_0p0003_mbp_mean}_{\pm\result{ext_code_llama_dora_0p0003_mbp_sd}}$ \\
 & \textcolor{PiSSAColor}{PiSSA} & $1\times10^{-4}$ & $\result{ext_code_llama_pissa_0p0001_he_mean}_{\pm\result{ext_code_llama_pissa_0p0001_he_sd}}$ & $\result{ext_code_llama_pissa_0p0001_hep_mean}_{\pm\result{ext_code_llama_pissa_0p0001_hep_sd}}$ & $\result{ext_code_llama_pissa_0p0001_mb_mean}_{\pm\result{ext_code_llama_pissa_0p0001_mb_sd}}$ & $\result{ext_code_llama_pissa_0p0001_mbp_mean}_{\pm\result{ext_code_llama_pissa_0p0001_mbp_sd}}$ \\
 & \textcolor{MiLoRAColor}{MiLoRA} & $3\times10^{-4}$ & $\result{ext_code_llama_milora_0p0003_he_mean}_{\pm\result{ext_code_llama_milora_0p0003_he_sd}}$ & $\result{ext_code_llama_milora_0p0003_hep_mean}_{\pm\result{ext_code_llama_milora_0p0003_hep_sd}}$ & $\result{ext_code_llama_milora_0p0003_mb_mean}_{\pm\result{ext_code_llama_milora_0p0003_mb_sd}}$ & $\result{ext_code_llama_milora_0p0003_mbp_mean}_{\pm\result{ext_code_llama_milora_0p0003_mbp_sd}}$ \\
 & \textcolor{HRAColor}{HRA} & $1\times10^{-3}$ & $\result{ext_code_llama_hra_0p001_he_mean}_{\pm\result{ext_code_llama_hra_0p001_he_sd}}$ & $\result{ext_code_llama_hra_0p001_hep_mean}_{\pm\result{ext_code_llama_hra_0p001_hep_sd}}$ & $\result{ext_code_llama_hra_0p001_mb_mean}_{\pm\result{ext_code_llama_hra_0p001_mb_sd}}$ & $\result{ext_code_llama_hra_0p001_mbp_mean}_{\pm\result{ext_code_llama_hra_0p001_mbp_sd}}$ \\
\addlinespace[3pt]
\bottomrule\end{tabular*}

\end{center}
\end{table}

\FloatBarrier
\paragraph{Task gains and retention costs.}
All six methods improve Llama's HumanEval+ and MBPP+ means over the base model, with rankings that differ across benchmarks and retention. PiSSA leads on HumanEval+, alongside higher retention costs. Together, these results extend the distinction between task performance and forgetting to coding, with the size of the task gain depending on the backbone.
